\exercisegroup

在所有§ 5的练习中， \(\mathfrak{C}\) 表示一种平等主义理论。

\begin{enumerate}
\def\labelenumi{\arabic{enumi}.}
\item
  ，以及关系 \(x = y\) 在 \(x\) 在 \(\mathfrak{C}\) .
\item
  设 R 是 G 中的一个关系，且设 x, y 为不同的字母。则这些关系 \((\exists x)(x = y \text{ 且 } R)\) , \((y|x)R\) 在 G 中等价。
\item
  设R和S是C中的关系，设T是C中的项，并设x和y是不同的字母。假设y不是T的常量，且x不出现在T中。设 \(C'\) 为将S附加到C的公理后所得的理论。若R在 \(C'\) 中关于x是函数性的，且如果 \((T|y)S\) 是C中的一个定理，则关系 \((T|y)R\) 在C中关于X是函数性的。
\item
  设 R 和 S 是 C 中的关系，并设 x 是一个不属于 C 的常量的字母。若 R 在 C 中关于 x 是函数性的，且 \(R \leftrightarrow S\) 是 C 中的一个定理，则 S 在 C 中关于 x 是函数性的。
\item
  设 \(\pmb{R}\) , \(\pmb{S}\) , \(\pmb{T}\) 是 \(\mathfrak{C}\) 中的关系，并令 \(\pmb{x}\) 是一个字母。如果 \(\pmb{R}\) 关于 \(\pmb{x}\) 在 \(\mathfrak{C}\) 中是函数性的，证明以下关系是 \(\mathfrak{C}\) 中的定理 :
\end{enumerate}

\[
(\text { 非   } (\exists x) (R \text {   且   } S)) \iff (\exists x) (R \text {   且   } (\text { 非   } S)),
\]

\((\exists x)(R \text{ 且 } (S \text{ 且 } T)) \Longleftrightarrow (\exists x)(R \text{ 且 } S) \text{ 且 } (\exists x)(R \text{ 且 } T)),\) \((\exists x)(R \text{ 且 } (S \text{ 或 } T)) \Longleftrightarrow ((\exists x)(R \text{ 且 } S) \text{ 或 } (\exists x)(R \text{ 且 } T)).\)

\begin{enumerate}
\def\labelenumi{\arabic{enumi}.}
\setcounter{enumi}{5}
\item
  证明该方案 \((\exists x)R \Longrightarrow R\) 在以下内容中提供隐式公理 \(\mathfrak{C}\) ，那么 \(x = y\) 是 \(\mathfrak{C}\) 中的定理 （参见练习1）。
\item
  证明该方案 \((\pmb{R} \Longleftrightarrow \pmb{S}) \Longrightarrow (\tau_{\pmb{x}}(\pmb{R}) = \tau_{\pmb{x}}(\pmb{S}))\) 在以下内容中提供隐式公理 \(\mathfrak{C}\) ，那么 \(x = y\) 是 \(\mathfrak{C}\) 中的定理 。（取 \(\pmb{R}\) 为关系 \(x = x\) , \(\pmb{S}\) 该关系 \(x = y\) ，然后代入 \(x\) 为了 \(y\) 在由此得到的公理中。)\footnote{我们将看到，这一理论在集合论（§4，第8节）中不存在具有至少两个元素的模型（即在该关系中 \((\exists x)(\exists y)(x \neq y)\) （这是真的）。}
\end{enumerate}

